% =========================================================
% Proof: Infinite Neighbourhood Characterization
% Source: notes/limit-points/notes-limit-points.tex
% Difficulty: Easy
% =========================================================

\subsection*{Infinite Neighbourhood Characterization}
\label{prf:inf-nbhd}

\begin{remark*}[Return]
$\leftarrow$ \hyperref[cor:inf-nbhd]{Back to Corollary in Notes}
\end{remark*}

\begin{proof}
\Claim $p$ is a limit point of $E$ if and only if every ball $B_r(p)$
contains infinitely many points of $E$.

\Given A metric space $(X,d)$, $E \subseteq X$, $p \in X$.

\Goal To prove the biconditional.

\Strategy
$(\Leftarrow)$: Infinitely many points implies at least one point
$\neq p$, so the limit point condition holds.

$(\Rightarrow)$: Suppose some ball $B_r(p)$ contains only finitely many
points of $E \setminus \{p\}$. Show this contradicts $p$ being a limit
point by shrinking the ball to exclude all of them.

% -------------------------------------------------------
% YOUR PROOF HERE
% -------------------------------------------------------

\end{proof}

\begin{remark*}[Proof shape]
The interesting direction is $(\Rightarrow)$: a finite set has a minimum
positive distance from $p$; shrink the ball below that minimum.
\end{remark*}

\begin{remark*}[Dependencies]
\begin{itemize}
  \item Definition~\ref{def:limit-point}: limit point via balls.
  \item Theorem~\ref{thm:seq-char-limit}: sequential characterization.
  \item Finite sets have a closest point: $\min\{d(p,x) : x \in F, x \neq p\} > 0$.
\end{itemize}
\end{remark*}
